% The results paragraph of the talk group with planted flaws, and their fixes: shared by the slides and by
% report-flaws.tex, the clickable version to present
% \BadReport: a results paragraph with planted flaws. \Flaw{n}{text} marks a flaw, \Good{text} something done right;
% both are plain text in the task and marked in the answer
\colorlet{GoodGreen}{green!40!black}
\newcommand\Flaw[2]{#2}
\newcommand\Good[1]{#1}
% how bad each flaw is: the color of its badge
\colorlet{SevCritical}{Warn}\colorlet{SevProblematic}{orange!85!black}\colorlet{SevFine}{gray}
\csdef{flawsev@1}{Problematic}\csdef{flawsev@2}{Problematic}\csdef{flawsev@3}{Problematic}\csdef{flawsev@4}{Critical}
\csdef{flawsev@5}{Critical}\csdef{flawsev@6}{Critical}\csdef{flawsev@7}{Fine}\csdef{flawsev@8}{Critical}
\newcommand\FlawBadge[1]{\tikz[baseline=(b.base)]\node[circle,fill=Sev\csuse{flawsev@#1},text=white,font=\tiny\sbseries,inner sep=.4pt,minimum size=2.5mm](b){#1};}
\newcommand\FlawShown[2]{\textcolor{Sev\csuse{flawsev@#1}}{\sbseries #2}\,\FlawBadge{#1}}
\newcommand\BadReport{%
   \textsb{Results.} We asked \Flaw1{12 friends from our lab} to solve \Good{the same five tasks with our tool and with the state of the art (CatPurr)}.
   With our tool they were \Flaw2{\qty{38}{\percent} faster} (\Flaw3{\(t\)-test, \(p = 0.04\)}).
   \Flaw4{We removed two participants whose times looked unusual.}
   \Good{We also measured errors, satisfaction, and learnability}, \Flaw5{but only speed was significant}.
   Satisfaction did not differ, \Flaw6{so both tools are equally pleasant}.
   The effect is strongest for experts (\Flaw7{\(p = 0.03\) for more than five years of experience}).
   \Flaw8{Our tool therefore makes developers more productive.}}
% \FixItem{n}{problem}{fix}: one flaw of the report, with its fix, on one line;
% the problem carries its own punctuation
\newcommand\FixItem[3]{\FlawBadge{#1}~\textsb{#2} \textcolor{darkgray}{#3}\par\vspace{.7mm}}
\newcommand\FixLevel[2]{\textcolor{Sev#1}{\sbseries\small #2}\par\vspace{.6mm}}
% \ReportFixes: the fixes, sorted by how bad the flaw is; \FixGroup{overlay}{items} shows a group of them,
% \FixThreats{text} the closing remark
\newcommand\FixGroup[2]{\onslide<#1->{#2}}
\newcommand\FixThreats[1]{\onslide<5->{#1}}
\newcommand\ReportFixes{%
      \FixGroup2{\FixLevel{Critical}{Critical}
      \FixItem4{Exclusion after a look:}{fix the rule first}
      \FixItem5{One of four outcomes:}{report all, correct with Holm}
      \FixItem6{No difference is not equality:}{show the confidence interval}
      \FixItem8{Productivity:}{say what you measured (speed)}}
      \FixGroup3{\vspace{1mm}\FixLevel{Problematic}{Problematic (hides)}
      \FixItem3{Which test?}{paired \(t\)-test with \(t\), df, exact \(p\)}
      \FixItem2{Relative only:}{seconds and a \qty{95}{\percent} confidence interval}
      \FixItem1{Friends:}{recruit outside the lab, say how}}
      \FixGroup4{\vspace{1mm}\FixLevel{Fine}{Could be fine}
      \FixItem7{Expert subgroup:}{exploratory, test it on new data}}
      \FixThreats{\vspace{1.5mm}\textcolor{gray}{Perfect reporting is not always possible: name the \textsb{threats to validity}.}}}
